%%%%%%%%%%%%%%%%%
% This is an sample CV template created using altacv.cls
% (v1.3, 10 May 2020) written by LianTze Lim (liantze@gmail.com). Now compiles with pdfLaTeX, XeLaTeX and LuaLaTeX.
% (v1.6.5b, 27 Jun 2023) forked by Nicolás Omar González Passerino (nicolas.passerino@gmail.com)
%
%% It may be distributed and/or modified under the
%% conditions of the LaTeX Project Public License, either version 1.3
%% of this license or (at your option) any later version.
%% The latest version of this license is in
%%    http://www.latex-project.org/lppl.txt
%% and version 1.3 or later is part of all distributions of LaTeX
%% version 2003/12/01 or later.
%%%%%%%%%%%%%%%%

%% If you need to pass whatever options to xcolor
\PassOptionsToPackage{dvipsnames}{xcolor}

%% If you are using \orcid or academicons
%% icons, make sure you have the academicons
%% option here, and compile with XeLaTeX
%% or LuaLaTeX.
% \documentclass[10pt,a4paper,academicons]{altacv}

%% Use the "normalphoto" option if you want a normal photo instead of cropped to a circle
% \documentclass[10pt,a4paper,normalphoto]{altacv}

%% Fork (before v1.6.5a): CV dark mode toggle enabler to use a inverted color palette.
%% Use the "darkmode" option if you want a color palette used to 
% \documentclass[10pt,a4paper,ragged2e,withhyper,darkmode]{altacv}

\documentclass[10pt,a4paper,ragged2e,withhyper]{altacv}

%% AltaCV uses the fontawesome5 and academicons fonts
%% and packages.
%% See http://texdoc.net/pkg/fontawesome5 and http://texdoc.net/pkg/academicons for full list of symbols. You MUST compile with XeLaTeX or LuaLaTeX if you want to use academicons.

% Change the page layout if you need to
\geometry{left=1.2cm,right=1.2cm,top=1cm,bottom=1cm,columnsep=0.75cm}

% The paracol package lets you typeset columns of text in parallel
\usepackage{paracol}

% Change the font if you want to, depending on whether
% you're using pdflatex or xelatex/lualatex
\ifxetexorluatex
  % If using xelatex or lualatex:
  \setmainfont{Roboto Slab}
  \setsansfont{Lato}
  \renewcommand{\familydefault}{\sfdefault}
\else
  % If using pdflatex:
  \usepackage[rm]{roboto}
  \usepackage[defaultsans]{lato}
  % \usepackage{sourcesanspro}
  \renewcommand{\familydefault}{\sfdefault}
\fi

% Fork (before v1.6.5a): Change the color codes to test your personal variant on any mode
\ifdarkmode%
  \definecolor{PrimaryColor}{HTML}{C69749}
  \definecolor{SecondaryColor}{HTML}{D49B54}
  \definecolor{ThirdColor}{HTML}{1877E8}
  \definecolor{BodyColor}{HTML}{ABABAB}
  \definecolor{EmphasisColor}{HTML}{ABABAB}
  \definecolor{BackgroundColor}{HTML}{191919}
\else%
  \definecolor{PrimaryColor}{HTML}{001F5A}
  \definecolor{SecondaryColor}{HTML}{0039AC}
  \definecolor{ThirdColor}{HTML}{F3890B}
  \definecolor{BodyColor}{HTML}{666666}
  \definecolor{EmphasisColor}{HTML}{2E2E2E}
  \definecolor{BackgroundColor}{HTML}{FFFFFF}
\fi%

\colorlet{name}{PrimaryColor}
\colorlet{tagline}{SecondaryColor}
\colorlet{heading}{PrimaryColor}
\colorlet{headingrule}{ThirdColor}
\colorlet{subheading}{SecondaryColor}
\colorlet{accent}{SecondaryColor}
\colorlet{emphasis}{EmphasisColor}
\colorlet{body}{BodyColor}
\pagecolor{BackgroundColor}

% Change some fonts, if necessary
\renewcommand{\namefont}{\Huge\rmfamily\bfseries}
\renewcommand{\personalinfofont}{\small\bfseries}
\renewcommand{\cvsectionfont}{\LARGE\rmfamily\bfseries}
\renewcommand{\cvsubsectionfont}{\large\bfseries}

% Change the bullets for itemize and rating marker
% for \cvskill if you want to
\renewcommand{\itemmarker}{{\small\textbullet}}
\renewcommand{\ratingmarker}{\faCircle}

%% Désactiver la césure automatique des mots en fin de ligne
\hyphenpenalty=10000
\exhyphenpenalty=10000
\tolerance=1000
\emergencystretch=3em

%% sample.bib contains your publications
%% \addbibresource{main.bib}

\begin{document}
    \name{Angenor N'Gouandi}
    \tagline{CHEF DE PROJET IT | PILOTAGE DE PROJETS \& GOUVERNANCE SI}
    %% You can add multiple photos on the left or right
    \photoL{4cm}{angenor}

    \personalinfo{
        \email{angenor99@gmail.com}
        \smallskip
        \phone{+225 05-45-29-28-02}
        \printinfo{\faBirthdayCake}{Né le 17/12/1999}
        \printinfo{\faFlag}{Nationalité ivoirienne}
        \printinfo{\faHome}{Domicilié à Abidjan, Côte d'Ivoire}
        %\\linkedin{angenor-ngouandi-463708186}
        \printinfo{\faGlobe}{Portfolio}[https://angenor.firebaseapp.com/]
        %\github{githubUser}
        %\npm{npmUser}
        %\dev{devtoUser}
        %\medium{nicolasomar}
        %% You MUST add the academicons option to \documentclass, then compile with LuaLaTeX or XeLaTeX, if you want to use \orcid or other academicons commands.
        % \orcid{0000-0000-0000-0000}
        %% You can add your own arbtrary detail with
        %% \printinfo{symbol}{detail}[optional hyperlink prefix]
        % \printinfo{\faPaw}{Hey ho!}[https://example.com/]
        %% Or you can declare your own field with
        %% \NewInfoFiled{fieldname}{symbol}[optional hyperlink prefix] and use it:
        % \NewInfoField{gitlab}{\faGitlab}[https://gitlab.com/]
        % \gitlab{your_id}
    }
    
    \makecvheader
    %% Depending on your tastes, you may want to make fonts of itemize environments slightly smaller
    % \AtBeginEnvironment{itemize}{\small}

    \begin{quote}
    \normalsize Chef de projet IT et ingénieur logiciel avec plus de 5 ans d'expérience dans le pilotage de projets de systèmes d'information : cadrage des besoins, planification, coordination des équipes et des parties prenantes, suivi de la recette et mise en production. Pratique des approches Agile/Scrum et du cycle en V (classique) ainsi que des référentiels de gouvernance et de gestion des services IT (COBIT, ITIL). Expert Informatique pour l'Organisation Internationale de la Francophonie, architecte technique pour l'Université Senghor et formateur technique, avec plus de 5 plateformes en production auprès d'institutions publiques et internationales. Master 2 en Système d'Information et Génie Logiciel.
    \end{quote}
    \vspace{-0.5em}

    %% Set the left/right column width ratio to 6:4.
    \columnratio{0.25}

    % Start a 2-column paracol. Both the left and right columns will automatically
    % break across pages if things get too long.
    \begin{paracol}{2}
        % ----- TECH STACK -----
        \cvsection{COMPÉTENCES}
            \textbf{Méthodologies \& Gouvernance} \\[0.3em]
            \cvtag{Agile / Scrum}
            \cvtag{Cycle en V / classique}
            \cvtag{ITIL}
            \cvtag{COBIT}
            \vspace{0.5em}

            \textbf{Gestion de projet} \\[0.3em]
            \cvtag{Cadrage \& planification}
            \cvtag{Pilotage \& coordination}
            \cvtag{Parties prenantes}
            \cvtag{Recette \& mise en production}
            \vspace{0.5em}

            \textbf{Architecture applicative} \\[0.3em]
            \cvtag{Architecture / SI}
            \cvtag{API REST}
            \cvtag{Intégration continue}
            \vspace{0.5em}

            \textbf{Développement} \\[0.3em]
            \cvtag{Python}
            \cvtag{Laravel}
            \cvtag{VueJs}
            \cvtag{Flask}
            \cvtag{Flutter}
            \vspace{0.5em}

            \textbf{Qualité / Tests} \\[0.3em]
            \cvtag{Squash-TM}
            \cvtag{Mantis}
            \cvtag{Selenium}
            \vspace{0.5em}

            \textbf{Bases de données} \\[0.3em]
            \cvtag{Supabase}
            \cvtag{Firebase}
            \cvtag{MySQL}
            \vspace{0.5em}

            \textbf{Outils / Infra} \\[0.3em]
            \cvtag{Git}
            \cvtag{Microsoft Azure}
            \cvtag{Linux}
        % ----- TECH STACK -----

        % ----- POINTS FORTS -----
        \cvsection{Points forts}
            \cvachievement{}{Pilotage multi-projets en production}{5 plateformes}
            \cvachievement{}{Gouvernance \& cycle de vie SI}{Cadrage, déploiement, recette, évolution}
            \cvachievement{}{Institutions publiques \& internationales}{OIF, ANSUT, ESATIC, Université Senghor}
            \cvachievement{}{Finaliste du Hackathon Francophone IA}{ESG Mefali — Top 10/200, IFDD \& Gouv. du Québec, 2026}
        % ----- POINTS FORTS -----

        % ----- LANGUAGES -----
        \cvsection{Langues}
            \cvlang{Français}{Langue maternelle} \\
            \cvlang{Anglais}{Intermédiaire}
        % ----- LANGUAGES -----

        % ----- permis -----
        \cvsection{Permis}
            \cvlang{Catégorie}{B, C, D, E} \\
        % ----- permis -----

        % ----- INTERETS -----
        \cvsection{Intérêts}
            \cvtag{Qualité logicielle}
            \cvtag{Pédagogie}
        % ----- INTERETS -----
        
        %% Switch to the right column. This will now automatically move to the second
        %% page if the content is too long.
        \switchcolumn
        
        
        
        % ----- EXPERIENCE -----
        \cvsection{Expériences}

            % \cvevent{Consultant développeur mobile}{Beyima}{20 Juillet 2022 -- Aujourd'hui}{Abidjan, Côte d'Ivoire}
            % \begin{itemize}
            %     \item Définition de l'architecture technique des applications mobiles.
            %     \item Prise de décisions sur les technologies et frameworks à utiliser.
            %     \item Participation active au développement des applications (iOS et/ou Android).
            %     \item Optimisation des performances et des expériences utilisateurs des applications.
            %     \item \textcolor{SecondaryColor}{\textbf{Outils: }} Flutter/Dart, Firebase, Laravel API
            % \end{itemize}
            % \divider
            
            
            
            \cvevent{Expert Informatique \& Chef de projet}{Organisation Internationale de la Francophonie - Institut de la Francophonie pour le Développement Durable}{Mai 2021 -- En cours(Fini en Déc. 2027)}{Québec, Canada (Télétravail)}
            \begin{itemize}
                \item Pilotage de la plateforme Epavillon Climatique — \cvreference{epavillonclimatique.francophonie.org}{https://epavillonclimatique.francophonie.org/} (gestion des activités du Pavillon de la Francophonie lors des COP pour les 53 États membres).
                \item Cadrage des besoins, définition de l'architecture technique et des choix technologiques de la plateforme.
                \item Conduite des phases de conception, développement et maintenance corrective et évolutive en approche itérative (Agile).
                \item Analyse des impacts techniques des nouvelles solutions et intégration d'API tierces (Zoom, YouTube, OpenRouter).
                \item Contrôle de conformité aux spécifications fonctionnelles et techniques avant chaque mise en production (gestion des changements, démarche ITIL).
                \item Pilotage de la recette et du suivi des anomalies avec les équipes métier ; conseil et accompagnement des parties prenantes.
                \item \textcolor{SecondaryColor}{\textbf{Outils: }} Python, Laravel, VueJs, Supabase, Flutter, API REST.
            \end{itemize}
            \divider

            \cvevent{Chef de projet (consultant) \& Développeur Full Stack}{Université Senghor — Université internationale de langue française, Opérateur direct de la Francophonie}{Décembre 2025 -- Décembre 2026}{Alexandrie, Égypte (Télétravail)}
            \begin{itemize}
                \item Site en production : \cvreference{usenghor-francophonie.org}{https://usenghor-francophonie.org/}.
                \item Conduite du projet de bout en bout en mode itératif (Agile) : recueil et analyse des besoins, cadrage, planification, conception, développement, tests et recette, jusqu'à la mise en production.
                \item Conception de l'architecture technique d'un site institutionnel trilingue (FR/EN/AR, RTL) : monorepo Nuxt 4 (SSR) + FastAPI, urbanisation en 12 services fonctionnels (\textasciitilde65 tables PostgreSQL).
                \item Définition et pilotage d'un modèle de gouvernance SI multi-acteurs pour l'administration décentralisée par les équipes des \textbf{14 pays} de représentation (RBAC, périmètres de contribution, journal d'audit).
                \item Coordination de la mise en production : Docker Compose, Nginx (reverse proxy, SSL/TLS, rate limiting), déploiement automatisé, sécurité JWT + refresh tokens, en-têtes HTTP sécurisés.
                \item Volumétrie livrée : 70+ pages, 135 composants Vue, 51 endpoints API, 15 services métier.
                \item \textcolor{SecondaryColor}{\textbf{Outils: }} Nuxt 4, Vue 3, Tailwind CSS 4, FastAPI, Python 3.14, SQLAlchemy 2 (async), Pydantic 2, PostgreSQL 16, Docker, Nginx.
            \end{itemize}
            \divider


            \cvevent{Enseignant vacataire en Informatique}{École Supérieure Africaine des TIC (ESATIC), Ministère de l'enseignement supérieur et de la recherche scientifique}{Octobre 2023 -- Aujourd'hui}{Abidjan, Côte d'Ivoire}
            \begin{itemize}
                \item Formation des étudiants de Licence 3 de la filière Développement d'Applications et Systèmes d'Information (DASI) en {\textbf{Programmation Web et Multimédia}}.
                \item Membre du jury des soutenances de Licence 3.
                \item \textcolor{SecondaryColor}{\textbf{Outils: }} HTML, CSS, PHP, Javascript, Modélisation UML
            \end{itemize}
            \divider

            \cvevent{Enseignant vacataire en sécurité Informatique}{Université Polytechnique de Bingerville}{Mai 2026 -- Aujourd'hui}{Abidjan, Côte d'Ivoire}
            \begin{itemize}
                \item Formation des étudiants de Master 1 de la filière ICIA en {\textbf{Développement d'applications sécurisé \& Audit de code}}.
                \item \textcolor{SecondaryColor}{\textbf{Outils: }} HTML, CSS, PHP, Javascript, UML, FastAPI
            \end{itemize}
            \divider

            \cvevent{Formateur technique ponctuel}{RMO Capital Humain \& Orange Côte d'Ivoire}{2023 -- 2025}{Abidjan, Côte d'Ivoire}
            \begin{itemize}
                \item Python pour des développeurs expérimentés chez RMO (RH, 30 000+ employés, 2025) ; Flutter pour des porteurs de projets chez Orange Côte d'Ivoire (2023).
            \end{itemize}

        \divider

        \cvevent{Développeur / Assurance Qualité}{Agence Nationale Du Service Universel Des Télécommunications (ANSUT)}{Avril 2019 -- Août 2019}{Abidjan, Côte d'Ivoire}
        \begin{itemize}
            \item Étude comparative de solutions open source de gestion des tests et de suivi des anomalies, puis choix, déploiement et configuration sur Microsoft Azure (Linux/Tomcat) de Squash-TM et Mantis Bug Tracker, pour outiller les chefs de projet de l'ANSUT : centralisation du suivi des projets sectoriels en remplacement des fichiers Excel.
            \item Outils destinés aux chefs de projet, interface entre la maîtrise d'ouvrage (MOA) et la maîtrise d'œuvre (MOE).
            \item Automatisation de tests via Python et Selenium IDE.
            \item \textcolor{SecondaryColor}{\textbf{Outils: }} Squash-TM, Mantis, Selenium IDE, Python, Azure, Tomcat, Linux.
        \end{itemize}

        % ----- PROJETS -----
        \cvsection{Projets notables}
            \cvevent{ESG Mefali — Conseiller ESG IA \textnormal{\textit{(Finaliste Hackathon Francophone IA)}}}{IFDD, Fondation Green Open Lab \& Gouvernement du Québec}{Mars 2026 -- En cours}{}
            \begin{itemize}
                \item Annonce officielle des 10 finalistes : \cvreference{linkedin.com/posts/fondation-green-openlab}{https://www.linkedin.com/posts/fondation-green-openlab_ia-francophonie-daezveloppementdurable-activity-7444395359969193984-NSg3}.
                \item Plateforme conversationnelle IA pour la finance durable des PME africaines francophones : diagnostic ESG, matching financement vert (GCF, FEM, BOAD, BAD), scoring de crédit alternatif, calcul d'empreinte carbone, extension Chrome de pré-remplissage de candidatures.
                \item \textcolor{SecondaryColor}{\textbf{Outils: }} Nuxt 4, FastAPI, Python, Claude API via OpenRouter, LangGraph/LangChain, PostgreSQL + pgvector.
            \end{itemize}
            \divider

            \cvevent{Observatoire des Mines de Madagascar}{\cvreference{miningobs.mg}{https://miningobs.mg/}}{Novembre 2025}{}
            \begin{itemize}
                \item Plateforme de l'Observatoire des Mines de Madagascar.
            \end{itemize}
            \divider

            \cvevent{Suivi des revenus miniers}{\cvreference{kaominina-mangarahara.mg}{https://kaominina-mangarahara.mg/}}{Novembre 2025}{}
            \begin{itemize}
                \item Plateforme numérique de suivi de l'utilisation des revenus miniers des collectivités territoriales de Madagascar.
            \end{itemize}
            \divider

            \cvevent{Le Carré des Études}{\cvreference{lecarredesetudes.com}{https://lecarredesetudes.com/}}{Mars 2026}{}
            \begin{itemize}
                \item Plateforme numérique du premier magazine dédié aux étudiants ivoiriens.
            \end{itemize}
            \divider

            \cvevent{Alumni ESATIC — Développeur front-end}{\cvreference{alumni-esatic.com}{https://alumni-esatic.com/}}{2025}{}
            \begin{itemize}
                \item Plateforme communautaire des anciens de l'ESATIC, développée en \textbf{équipe Agile/Scrum} (chef de projet, dev back-end, dev front-end, DevOps) : réunions régulières, backlog, répartition des tâches.
            \end{itemize}


        % ----- MISSIONS -----
        \cvsection{MISSIONS INTERNATIONALES}
            \cvevent{Conférences des Parties (COP 26 à 30)}{Organisation Internationale de la Francophonie}{2021 -- 2025}{}
            \begin{itemize}
                \item Support technique du Pavillon de la Francophonie lors des COP : Glasgow, Royaume-Uni (2021) — Charm el-Cheikh, Égypte (2022) — Dubaï, Émirats arabes unis (2023) — Bakou, Azerbaïdjan (2024) — Belém, Brésil (2025).
            \end{itemize}
            \divider

            \cvevent{Salon des sciences, technologies et innovations environnementales}{Organisation Internationale de la Francophonie}{2024 -- 2025}{}
            \begin{itemize}
                \item Support technique: Yaoundé, Cameroun (Février 2025) — Kinshasa, RDC (Novembre 2024).
            \end{itemize}
        % ----- MISSIONS -----
        
        % ----- EDUCATION -----
        \cvsection{FORMATIONS}
            \cvevent{Master 2 Système d'Information et Génie Logiciel}{École Supérieure Africaine des TIC (ESATIC)}{2020 -- 23/09/2021}{Abidjan, Côte d'Ivoire}
            \divider

            \cvevent{Licence en Système Réseaux Informatiques et Télécommunication}{École Supérieure Africaine des TIC (ESATIC)}{2018 -- 20/09/2019}{Abidjan, Côte d'Ivoire}
            \divider

            \cvevent{Baccalauréat série C}{Lycée Classique d'Abidjan}{2016 -- 25/07/2016}{Abidjan, Côte d'Ivoire}

        % ----- EDUCATION -----
        
    \end{paracol}
\end{document}